\section{序論と主結果}

本稿は、重み付き射影三次元スタック
\[
 X_m=\mathbb P_{\mathrm{stack}}(1,1,1,m),\qquad m\ge3
\]
のPoisson構造のパラメータ空間と、その明示的な非可換量子化を調べる。
大域的二重ベクトル場のアフィン空間にJacobi方程式を課したスキームを
$\mathscr P_m$ と書く。スカラー倍と重み付き座標変換は商せず、
方程式のイデアルを根基化しない。
この規約の下では、被約な族が一つの座標変換軌道で稠密に覆われる場合にも、
パラメータ空間の一般点に非自明な冪零方向が残る。

主結果は、接する直線と二次曲線から生じる共鳴系列の全形式近傍と、
その一般点を量子化する通常の4-Calabi--Yau代数の中心を決定することである。
幾何側には任意に大きい一般点の長さが現れ、
量子代数側には奇数重みの端でDu Val曲面が現れる。
これらを、別の引戻し成分の厚み、および重み三の成分・境界の計算と合わせて記述する。

\subsection{接線共鳴の成分系列と全形式近傍}

$S_d=\mathbb C[x,y,z]_d$、$N_m=\binom{m+2}{2}$ と置く。
Poisson構造の標準形には
$c\in S_3,D\in\mathfrak{sl}_3,[F]\in S_{m+1}^{\oplus3}/\mathbf xS_m$
を用いる。その全Jacobi式は
\[
 D(c)=0,\qquad
 \frac{m+3}{3}F(c)-\frac m3(\mathbf x\times D)\cdot F-c\operatorname{div}F=0
\]
である。

\begin{theorem}[接線共鳴系列]
\label{thm:main-tangent-v34}
$0\le j\le\lfloor(m-1)/2\rfloor$、$r=m-2j$ とし、
\[
 q=xz-y^2,\quad c=xq,\quad
 D=-\frac{4m+3-6j}{2m}(x\partial_y+2y\partial_z),\quad
 F_j=q^j(0,yz^r/2,z^{r+1})
\]
と置く。組 $(c,D)$ の $\operatorname{SL}_3$ 軌道上で全混合核を動かした族の閉包
$\mathscr E_{m,j}$ は、次元 $N_m+8$ の既約成分である。
その一般点の局所環は $K[\gamma]/(\gamma^{j+1})$ である。
従って固定した $m$ で、一般点の長さが
$1,2,\ldots,\lceil m/2\rceil$ となる相異なる成分を得る。

$F=0$ の基準点では、全完備局所環は
\[
 \boxed{
 \mathbb C[[s_1,\ldots,s_{N_m+7},\kappa,\gamma,u]]/
              (\kappa u,\gamma^{j+1}u)
 }
\]
である。大きい方の枝は次元 $N_m+9$ の $L(1,2,m)$ 成分で、
小さい方が $\mathscr E_{m,j}$ である。
小さい方の被約成分は重み付き座標変換群の稠密な一軌道を持つ。
\end{theorem}

証明では、基底 $D(c)=0$ の七軌道方向を除く二変数切片を取り、
実際の $w$ 移動の核を非被約基底上でも自由な直和因子として除く。
残る一列のSchur補完の可能な次数は四つの余核重量から制約される。
直線と二次曲線を横断させた場合の全混合核を多項式の次数降下で解き、
二つの特異点での評価を使うと、最初の非零障害が正確に $j+1$ 次であると分かる。
有限個の重みの階数検算から一般公式を外挿する証明ではない。

\subsection{Calabi--Yau量子化とDu Val中心}

今度は $j,r\ge0$、$m=2j+r\ge3$ とし、$r=0$ も含める。
導分
\[
 N=x\partial_y+2y\partial_z-q^jz^r\partial_w,\qquad
 X_\alpha=\alpha x\partial_x+(\alpha+1)y\partial_y+(\alpha+2)z\partial_z
                +(m\alpha+m+r+1)w\partial_w
\]
は $[N,X_\alpha]=N$ を満たす。
Coll--Gerstenhaber--Giaquintoの有限和
\[
 a\star b=\sum_{n\ge0}s^nN^n(a)\binom{X_\alpha}{n}(b)
\]
から得る重み $(1,1,1,m)$ の代数を $\mathcal S_{j,r}(\alpha,s)$ とする。

\begin{theorem}[Nakayama公式と中心]
\label{thm:main-quantum-v34}
$\mathcal S_{j,r}(\alpha,s)$ はnoetherian整域かつAS正則四次元代数で、
Hilbert級数は $((1-t)^3(1-t^m))^{-1}$、Gorenstein parameterは $m+3$ である。
$T=1+sN$ と置くと、Nakayama自己同型は
\[
 \boxed{\mu=T^{(m+3)\alpha+m+r+3}}
\]
となる。従って $s\ne0$ のとき通常の4-Calabi--Yau条件は
$\alpha=-(m+r+3)/(m+3)$ と同値である。
このCalabi--Yau族の中心は
\[
 \begin{array}{c|c}
 r\ge2&\mathbb C\\
 r=0&\mathbb C[q,J]\\
 r=1&\mathbb C[S,P,R]/(R^2-SP^2-\tfrac{16}{81}S^{m+3})
 \end{array}
\]
であり、最後の場合は $D_{m+4}$ 型のDu Val曲面である。
ここで $J,S,P,R$ は本文で明示する多項式であり、中心の生成性も証明する。
\end{theorem}

三変数部分代数はLecoutre--Sierraの $R(2,\alpha)$ である。
その既知のNakayama公式\cite[Theorem 4.16]{LecoutreSierra2019}とOre拡大の制限公式を用い、
最後の重み $m$ の生成元に関する剛性を証明することで、上の四変数の公式を得る。
中心の決定には、局所冪零導分の核を四生成元一関係で完全に記述し、
その零重量部分を求める。多項式不変環の生成に関する方法は
Maubach\cite{Maubach2000}の研究に位置付けられる。

$r>0$ のCalabi--Yau族は、前定理の成分の稠密な軌道を実際に量子化する。
全ての族で標準 $(m+3)$ 項列はfull strongであり、
$(0,1,m,m+1)$ の四項部分列は許容部分圏を与える。
特に $m=3$ では、二つの十八次元成分の双方が非可換VK圏の標準六項実現を持つ。
固定した非零量子代数自体の変形空間の一般点長さを、前定理で計算したという意味ではない。

二つの定理を奇数重みの最後の共鳴に適用すると、次の対応を得る。
\[
 \boxed{
 m=2\ell-1,\ j=\ell-1,\ r=1\quad(\ell\ge2):\qquad
 \text{一般点の長さ }\ell,\qquad
 \text{量子中心の型 }D_{2\ell+3}.
 }
\]
一方、次数零の量子関数斜体は全て
$\mathbb C(C)$ 上の第一Weyl斜体であり、$C=J^2/q^{m+r+1}$ である。
従って多項式中心の差は、この双有理表示だけでは検出できない。
異なる重みの射影圏は $K_0$ の階数 $m+3$ が異なるため導来同値でないが、
同じ重みの異なる共鳴の射影圏の比較は残る。

\subsection{引戻し成分に現れる別の厚みと圏の変形基底}

\begin{theorem}[引戻し成分]
\label{thm:main-v33}
全ての $m\ge3$ で、$\mathbb P(1,1,m)$ からの引戻し族の閉包は
次元 $N_m+9$ の既約成分であり、一般点の局所環は
\[
 K[a,b]/(a^{m+2},ab,b^2)
\]
となる。その長さは $m+3$ である。
$L(1,1,1,m)$ 成分との一般的な交叉では、全完備局所環は
\[
 \mathbb C[[s_1,\ldots,s_{N_m+8},a,b,u]]/(ab,b^2,a^{m+2}u)
\]
となる。この局所環はCohen--Macaulayであるが完全交叉ではない。
\end{theorem}

接線系列の一変数のArtin環と異なり、この引戻し成分の一般点には
二つの独立な接方向が残る。
全重みでの証明には、実際の座標移動核、重量による障害次数の制約、
三角形への変形と三交点での評価を用いる。
この族には発散零の微分Ore拡大による、別の明示的な4-Calabi--Yau量子化も与える。

標準例外列代数
\[
 \Lambda_m=\operatorname{End}_{X_m}
       \left(\bigoplus_{i=0}^{m+2}\mathcal O_{X_m}(i)\right)
\]
のminiversal変形基底は、非標準的な形式座標の選択の下で
$\mathscr P_m$ の原点完備化に同型である。
変数数は $m^2+6m+26$、独立な二次生成元の数は $10+\binom{m+5}{2}$ となる。
斉次錐の一般点長さは原点完備化でも保存されるので、
上の二種類の非被約構造はこの圏の変形基底にも現れる。
これは自己同型を含む変形groupoidのfine表現を主張するものではない。

\subsection{重み三の成分と境界}

重み三の問題は、Veronese--Kronecker四項切り出しを
標準六項へ補完する問題から生じる\cite{VKAmbient32}。
以下の相異なる十成分を認定する。
\begin{center}\small
\begin{tabular}{lcl}\toprule
一般的な構成&次元&一般点の性質\\\midrule
$\operatorname{PB}(\mathbb P^2)$&35&滑らか\\
有理型 $R(1,5)$&28&滑らか\\
対数型 $L(1,1,4)$&23&滑らか\\
有理型 $R(2,4)$&22&滑らか\\
有理型 $R(3,3)$&19&滑らか\\
対数型 $L(1,2,3)$&19&滑らか\\
対数型 $L(1,1,1,3)$&19&滑らか\\
$\operatorname{PB}(\mathbb P(1,1,3))$&19&長さ六\\
接線共鳴 $a=-5/2$&18&滑らか（長さ一）\\
接線共鳴 $a=-3/2$&18&長さ二\\\bottomrule
\end{tabular}
\end{center}
$R(p,q)$ は $pf\,dg-qg\,df$ の族、$L(d_1,\ldots,d_k)$ は
$(\prod f_i)\sum\lambda_i\,df_i/f_i$ の族の閉包を表す。
十成分が全てであること、三十五が最大次元であることは主張しない。

接線型の四つの非零共鳴では、基底変数を含む全形式近傍を
有理係数のSchur消去でも独立に求める。
非共点三直線の特殊点では五つの局所枝を持つ非被約環を決定し、
$D=0$ の全非零三次式の混合核は留数で記述する。
$c=0$ の閉部分の被約化は次元 $35,20,18$ の三成分からなり、
全て既知の大域成分に含まれることも示す。

\subsection{先行研究との関係と論証の範囲}

三次元のPoisson構造と反標準束で捻った可積分一形式の対応、
有理型・対数型・引戻しという構成自体は古典的である。
通常の射影三空間におけるPoisson構造と量子代数の関係はPym\cite{Pym}、
関連する変形論は\cite{Molinuevo2016}、
閉有理一形式の留数は\cite{PereiraClosedForms}を参照されたい。
CGG公式、AS正則性に関するOre拡大定理、既知の三変数Nakayama公式を
本稿の新しい構成原理として扱わない。
本稿で直接求めるのは、この重み付き族の成分、全局所環、
非線形な最後の生成元の剛性、および中心の具体的な表示である。

重み付き射影空間を標的とする引戻しの安定性\cite{GAMQV2026}は、
通常の射影空間を始域とする場合などに仮定を課しており、
その一般被約性を本稿の重み付き三次元始域へ直接適用することはできない。
葉層スキームの非被約性自体にも先行例がある
\cite[Remark 4.2]{BKM2026}。
本稿では一般点のArtin環と交叉環を決定し、整数パラメータとその長さの関係まで証明する。
一般重みの成分定理・中心の生成性と、有限の有理数検証の役割は明確に区別する。
